
\exo

Développer, réduire et ordonner les expressions suivantes :

\vspace{-6pt}
\setlength{\columnsep}{1cm}
\setlength{\columnseprule}{0pt}
\begin{multicols}{2}
\begin{enumerate}
\item $A=(4x-5)(2x-3)$.
\smallskip
\item $B=-3(2-6x)(6-2x)$.
\smallskip
\item $C=2(x+3)(5x+1)$.
\smallskip
\item $D=(3x-2)(5-x)-4x(x+6)$.
\smallskip
\item $E=(x+5)^2-21+2x$.
\smallskip
\item $F=4(2x-3)^2$.

\item $G=3(t-2)^2+1$.
\smallskip
\item $H=-2(x+4)^2+4x+7$.
\smallskip
\item $I=(x+7)^2+2x+4$.
\smallskip
\item $J=-3(x-4)^2+11$.
\smallskip
\item $K=(2x-5)(2x+5)-(3x+5)^2$.
\smallskip
\item $L=(x-1)^2\times (x+2)$.
\end{enumerate}
\end{multicols}
\vspace{-12pt}

\exo

Montrer que les trois expressions $A$, $B$ et $C$ sont égales pour tout réel $t$ :
\[A=(2t-4)^2+12
\qpvq
B=4(t-2)^2+12
\qpvq
C=4(t-3)(t-1)+16.\]

\exo

Écrire les expressions suivantes sous la forme $a+b\sqrt{2}$, avec $a$ et $b$ entiers :

\vspace{-6pt}
\setlength{\columnsep}{1cm}
\setlength{\columnseprule}{0pt}
\begin{multicols}{2}
\begin{enumerate}
\item $A=(1+\sqrt{2})^2$.
\smallskip
\item $B=(3-\sqrt{2})^2$.

\item $C=(4+2\sqrt{2})^2$.
\smallskip
\item $D=(4+\sqrt{2})(4-\sqrt{2})$.
\end{enumerate}
\end{multicols}
\vspace{-12pt}

\exo

Factoriser les expressions suivantes en remarquant un facteur commun :

\vspace{-6pt}
\setlength{\columnsep}{1cm}
\setlength{\columnseprule}{0pt}
\begin{multicols}{2}
\begin{enumerate}
\item $A=(23x+1)(-17x+1)+(23x+1)^2$.
\smallskip
\item $B=(13x-4)(25x-11)-(13x-4)^2$.
\smallskip
\item $C=(8-18x)^2-(16x-3)(8-18x)$.

\item $D=(2x+3)(2x+4)-3(2x+3)$.
\smallskip
\item $E=(2x+3)(2x+4)+(2x+3)$.
\smallskip
\item $F=(2x-3)(4x+5)+(3-2x)$.
\end{enumerate}
\end{multicols}
\vspace{-12pt}

\exo

Factoriser les expressions suivantes en utilisant une identité remarquable :

\vspace{-6pt}
\setlength{\columnsep}{1cm}
\setlength{\columnseprule}{0pt}
\begin{multicols}{2}
\begin{enumerate}
\item $A=x^2-5$.
\smallskip
\item $B=(x+4)^2-5$.

\item $C=25-(2-x)^2$.
\smallskip
\item $D=(x+3)^2-(2x+1)^2$.
\end{enumerate}
\end{multicols}
\vspace{-12pt}

\exo

Factoriser les expressions suivantes :

\vspace{-6pt}
\setlength{\columnsep}{1cm}
\setlength{\columnseprule}{0pt}
\begin{multicols}{2}
\begin{enumerate}
\item $A=7x^3+28x^2$.
\smallskip
\item $B=4x^2-49$.
\smallskip
\item $C=2x^2-4$.
\smallskip
\item $D=\frac{4}{9}x^2-\frac{25}{49}$.

\item $E=x^2+2x+1$.
\smallskip
\item $F=9x^2+24x+16$.
\smallskip
\item $G=\frac{1}{3}x^2+2x+3$.
\end{enumerate}
\end{multicols}
\vspace{-12pt}

\exo

Simplifier, si possible, les quotients suivants :

\vspace{-6pt}
\setlength{\columnsep}{1cm}
\setlength{\columnseprule}{0pt}
\begin{multicols}{2}
\begin{enumerate}
\item $A=\dfrac{2x+4}{x+2}$.
\smallskip
\item $B=\dfrac{6x-4}{10x+20}$.

\item $C=\dfrac{5x^2+4x}{x}$.
\smallskip
\item $D=\dfrac{2x^2+3x}{x^2+x}$.
\end{enumerate}
\end{multicols}
\vspace{-12pt}

\exo

Écrire sous la forme d'une seule fraction, la plus simple possible :

\vspace{-6pt}
\setlength{\columnsep}{1cm}
\setlength{\columnseprule}{0pt}
\begin{multicols}{2}
\begin{enumerate}
\item $A=\dfrac{1}{x+1}-\dfrac{3}{x}$.
\smallskip
\item $B=\dfrac{2x+4}{x-2}+\dfrac{1}{2}$.

\item $C=\dfrac{4}{x-4}-\dfrac{3}{x+1}$.
\smallskip
\item $D=\dfrac{2x+2}{2x-1}+\dfrac{3x}{x+3}$.
\end{enumerate}
\end{multicols}
\vspace{-12pt}

\exo

Résoudre dans $\R$ les équations suivantes :

\vspace{-6pt}
\setlength{\columnsep}{1cm}
\setlength{\columnseprule}{0pt}
\begin{multicols}{3}
\begin{enumerate}
\item $(x+4)(x-7)=0$.
\smallskip
\item $(2x+3)(4x-5)=0$.

\item $x(5x-4)=0$.
\smallskip
\item $(15x-3)(3x+9)=0$.

\item $(2x-3)^2=0$.
\smallskip
\item $3x(x-5)=0$.
\end{enumerate}
\end{multicols}
\vspace{-12pt}

\exo

Résoudre dans $\R$ les équations suivantes :

\vspace{-6pt}
\setlength{\columnsep}{1cm}
\setlength{\columnseprule}{0pt}
\begin{multicols}{2}
\begin{enumerate}
\item $5x^2-6x=0$.
\smallskip
\item $(2x+1)(x+4)+(x+4)(3-5x)=0$.
\smallskip
\item $(x-7)(3x-5)-(9x-4)(x-7)=0$.

\item $4x^2+8x+4=0$.
\smallskip
\item $x(2x-3)=x(x+7)$.
\smallskip
\item $(4x-3)(x+5)=(4x-3)(2x-1)$.
\end{enumerate}
\end{multicols}
\vspace{-12pt}

\exo

\begin{enumerate}
\item Factoriser $x^2-16$, puis résoudre dans $\R$ l'équation $x^2-16=0$.
\item Résoudre dans $\R$ les équations suivantes :
\vspace{-6pt}
\setlength{\columnsep}{1cm}
\setlength{\columnseprule}{0pt}
\begin{multicols}{3}
\begin{enumerate}
\item $x^2-81=0$.
\smallskip
\item $x^2=15$.

\item $x^2+7=0$.
\smallskip
\item $3x^2=48$.

\item $2x^2-20=0$.
\smallskip
\item $4x^2-2=1$.
\end{enumerate}
\end{multicols}
\vspace{-12pt}
\end{enumerate}

\exo

Résoudre dans $\R$ les équations suivantes :

\vspace{-6pt}
\setlength{\columnsep}{1cm}
\setlength{\columnseprule}{0pt}
\begin{multicols}{4}
\begin{enumerate}
\item $\dfrac{x-2}{x+9}=0$.

\item $\dfrac{2x-7}{x+3}=0$.

\item $\dfrac{20-4x}{x-5}=0$.

\item $\dfrac{5x-1}{2x+3}=0$.
\end{enumerate}
\end{multicols}
\vspace{-12pt}

\exo

\vspace{-6pt}
\setlength{\columnsep}{1cm}
\setlength{\columnseprule}{0pt}
\begin{multicols}{2}
Une boîte avec couvercle a pour base un carré de côté $x$ centimètres et pour hauteur $2$ centimètres.
\begin{center}
\def\xmin{-0.5}
\def\xmax{4.6}
\def\ymin{-0.5}
\def\ymax{1.6}
\def\xunite{1.0cm}
\def\yunite{1.0cm}
\psset{unit=\xunite, xunit=\xunite, yunit=\yunite, algebraic=true, dimen=middle, dotstyle=*, dotsize=3pt 0, linewidth=0.8pt, arrowsize=3pt 2, arrowinset=0.25, comma=true}
\begin{pspicture*}(\xmin,\ymin)(\xmax,\ymax)
%\psgrid[xunit=\xunite, yunit=\yunite, subgriddiv=0, gridlabels=0, gridcolor=lightgray](0,0)(\xmin,\ymin)(\xmax,\ymax)
%\psaxes[labelFontSize=\scriptstyle, xAxis=true, yAxis=true, Dx=1, Dy=1, ticksize=-2pt 0, subticks=1]{->}(0,0)(\xmin,\ymin)(\xmax,\ymax)[{\footnotesize $x$},135][{\footnotesize $y$},-45]
\pspolygon[fillcolor=gray,fillstyle=solid,opacity=0.1](0,0)(3,0)(3,1)(0,1)
\pspolygon[fillcolor=gray,fillstyle=solid,opacity=0.1](3,0)(4.5,0.5)(4.5,1.5)(3,1)
\pspolygon[fillcolor=gray,fillstyle=solid,opacity=0.1](0,0)(1.5,0.5)(1.5,1.5)(0,1)
\pspolygon[fillcolor=gray,fillstyle=solid,opacity=0.1](1.5,0.5)(4.5,0.5)(4.5,1.5)
(1.5,1.5)
\pspolygon[fillcolor=gray,fillstyle=solid,opacity=0.1](0,0)(3,0)(4.5,0.5)(1.5,0.5)
\pspolygon[fillcolor=gray,fillstyle=solid,opacity=0.1](0,1)(3,1)(4.5,1.5)(1.5,1.5)
\begin{footnotesize}
\uput{5pt}[270](1.5,0){$x$}
\uput{5pt}[330](3.75,0.25){$x$}
\uput{5pt}[180](0,0.5){$2$}
\end{footnotesize}
\end{pspicture*}
\end{center}
\columnbreak
\begin{enumerate}
\item Montrer que l'aire de la surface extérieure de cette boîte est donnée, en \si{\cm^2}, par :
\[A(x)=2(x+2)^2-8.\]
\item Déterminer les valeurs de $x$ pour lesquelles cette aire est égale à \SI{90}{\cm^2}.
\end{enumerate}
\end{multicols}

\exo

Une expérience consiste à observer le développement de bactéries dans une boîte de Petri pendant $10$ jours. Le nombre de bactéries (en milliers) est donné en fonction du temps (en jours) par :
\begin{center}
$N(t)=(0,5t+1)^2$, pour tout nombre réel $t$ appartenant à $\intff{0}{10}$.
\end{center}
\begin{enumerate}
\item Donner une estimation du nombre de bactéries après un jour, puis après deux jours et demi.
\item Après quelle durée le nombre de bactéries a-t-il atteint \np{16000} ? Et \np{22090} ?
\end{enumerate}

\exo

La fonction $f$ est définie sur $\R$ par $f(x)=2x^2+4x-6$.
\begin{enumerate}
\item Vérifier que, pour tout réel $x$, l'expression $f(x)$ peut aussi s'écrire :
\[f(x)=2(x+1)^2-8\qetq f(x)=2(x-1)(x+3).\]
\item On dispose ainsi de trois formes de $f(x)$. En utilisant la forme la plus adaptée, déterminer :
\begin{enumerate}
\item les images de $0$, de $1$ et de $\sqrt{3}-1$ ;
\item les éventuels antécédents de $0$ et ceux de $-6$ ;
\item les abscisses des points de la courbe $\cC_f$ ayant pour ordonnée $24$.
\end{enumerate}
\end{enumerate}

